\scanpage{33}
trouver $U$ \inserted{comme dessus}, mais on peut trouver $U$ de
façon que l'on ait
\[
S_Y(U(a))=na
\]
d'où, posant $Z'(a)=Z(U(a))$, la relation
\[
Z'(S_Y(\unsure{\mathfrak{a}|\mathfrak{z}|\mathfrak{t}}))=nZ(\mathfrak{a})\qquad \mathfrak{a}\ \text{de degré 0 sur }Y.
\]
\deleted{Car on peut trouver $nS(\mathfrak{a})=U(a)$ si $a=S_Y(\mathfrak{a})$,
$S(\mathfrak{a})$ de la forme $U(\unsure{a|a'})$, car si $a'=S_Y(\mathfrak{a})$,
on aura $S_Y(U(a'))=na'=S_Y(n\mathfrak{a})$ donc
$U(a')=n\mathfrak{a}\unsure{)}$, et comme on aura $a'=na$ pour
$a$ convenable, il s'ensuivra $U(a)$}

Pour le voir, soit $a=S_Y(\mathfrak{a})$, on aura
$\deleted{\unsure{Z'(S}}S_Y(U(a))=na=S_Y(n\mathfrak{a})$ d'où \deleted{$U(a)=n\mathfrak{a}$,}
\[
U(a)=n\mathfrak{a}
\]
d'autre part on aura
\[
\begin{aligned}
\deleted{\unsure{S}}Z'(S_Y(U(a)))&=nZ(U(a))\\
\shortparallel\quad &\phantom{=}\shortparallel\\
Z'(na)\quad &\phantom{=}nZ'(a)\quad .
\end{aligned}
\]
\deleted{d'où, comme $U(a)=n\mathfrak{a}$}
Or tout \unsure{élément|élém} de $A_{0\,\mathrm{alb}}(Y)$ de degré 0 est
\deleted{$nZ'(S_Y(\mathfrak{a}))=n^2Z(\mathfrak{a})$}
de la forme $n\unsure{\mathscr L|\mathfrak b|L|\mathcal L}$ (car $A_Y(k)$ est divisible !)
donc appliquant ce qui précède à $\mathscr L$ au lieu de $\mathfrak{a}$,
on trouve
\[
Z'(S_Y(\mathfrak{a}))=nZ(\mathfrak{a}).
\]

On prend comme isogénie la multiplication par $n$.
% archive sheet number: 32
